\documentclass[12pt]{article}
\usepackage[utf8]{inputenc}
\usepackage[T1]{fontenc}
\usepackage[margin=1in]{geometry}
\usepackage{setspace}
\usepackage{xcolor}
\usepackage[colorlinks=true, urlcolor=blue!60!black, linkcolor=blue!60!black]{hyperref}
\setlength{\parindent}{0em}
\setlength{\parskip}{0.6em}
\renewcommand{\baselinestretch}{1.0}
\begin{document}
\begin{center}
{\LARGE \textit{Bradwell v. Illinois}, 83 U.S. 130 (1873)\par}
\vspace{0.5em}
{\large Case Brief --- Two-Page Edition\par}
\end{center}

\section{Memory Jogger}
A day after \textit{Slaughter-House}, the Court held 8--1 that Illinois could bar Myra Bradwell from practicing law because state bar admission is not a privilege of national citizenship, while Justice Bradley's concurrence invoked ``the law of the Creator'' to confine women to the domestic sphere (\href{https://supreme.justia.com/cases/federal/us/83/130/}{83 U.S. 130}).

\section{Detailed Case Facts}
Myra Bradwell, Vermont-born editor of the \textit{Chicago Legal News} and a longtime Chicago resident, applied to the Illinois Supreme Court for a law license with certificates of character and qualification. The court refused her, first on the ground of coverture and then because the legislature, adopting the English common law, never contemplated admitting women. She claimed a right to a license under Article IV, \S{} 2 and the Fourteenth Amendment's Privileges or Immunities Clause.

\section{Procedural History}
Bradwell sued out a writ of error to the U.S. Supreme Court, which affirmed the Illinois Supreme Court on April 15, 1873.

\section{Judicial Votes}
Affirmed, 8--1. Miller wrote for Clifford, Davis, Strong, and Hunt; Bradley, joined by Swayne and Field, concurred in the judgment; Chief Justice Chase, gravely ill, dissented without opinion.

\section{Holding}
Article IV, \S{} 2 does not protect a citizen against her own state's laws, and Bradwell was a citizen of Illinois. Admission to a state bar is not a privilege or immunity of United States citizenship, so states may set bar qualifications free of Fourteenth Amendment control. Bradley's views on women were never the Court's holding.

\section{Analysis of Opinions}
Miller's brief, doctrinal opinion avoided gender stereotypes: bar admission had never depended on U.S. citizenship, and \textit{Slaughter-House} made the regulation of law licenses a power never transferred to the federal government. Bradley, who had dissented in \textit{Slaughter-House} to protect the right to pursue a lawful occupation, reconciled his votes by asserting that this right had never been a fundamental privilege of women, citing coverture, ``nature herself,'' and the need for ``decision and firmness'' in the law. Chase dissented ``from the judgment of the Court and from all the opinions.''

\section{Examples: Future Applications}
\textbf{Same-Side:} a resident's Privileges or Immunities challenge to a 1,500-hour hair-braiding license fails because a state-regulated occupation is not a privilege of national citizenship. \textbf{Opposite-Side:} a state nursing school's exclusion of applicants by sex falls under modern Equal Protection doctrine, as in \textit{Mississippi University for Women v. Hogan}. \textbf{Fence-Sitter:} a federal district court's bar exclusion, given Miller's dicta that federal-court admission ``would relate to citizenship of the United States.''

\section{Critique}
Feminist scholars see \textit{Bradwell} as constitutionalized patriarchy that left women without federal recourse until \textit{Reed v. Reed} (1971). Originalists such as Timothy Sandefur argue that \textit{Slaughter-House} and \textit{Bradwell} mutilated the Clause's original meaning, though they divide over whether that meaning would have helped Bradwell; textualists fault Bradley for natural-law assumptions.

\section{Key Quotations}
\begin{quote}
``Man is, or should be, woman's protector and defender. The natural and proper timidity and delicacy which belongs to the female sex evidently unfits it for many of the occupations of civil life.''
\end{quote}
\noindent --- Justice Bradley, concurring in the judgment, 83 U.S. at 141.

\end{document}
